% \iffalse meta-comment
%
% hit-keylist.dtx 是各个类的键清单：这个类接受哪些信息字段与常量
%
% \fi
%
% \section{键清单（keylist）}
%
% \changes{v3.2a}{2026/08/10}{两个类的键清单合成一个文件，同名的那批提到
%   \texttt{shared-keylist} 守卫共用一份。原先 96 个名字在两个文件里各写一遍，
%   加一个公共字段要改两处，漏一处没有任何东西会说}
% 这个文件分三个守卫：\texttt{shared-keylist} 两个类都取，\texttt{thesis-keylist} 与
% \texttt{report-keylist} 各取各的。\file{hithesis.ins} 里一句
% \cs{from}\marg{hit-keylist.dtx}\marg{shared-keylist,thesis-keylist} 按文件顺序取出
% 共用那半加本类那半。
%
% 这一层只声明，不给值。值在 \file{hit-defaults.dtx}，机制在
% \file{hit-key-engine.dtx}，规矩见 CONTRIBUTING §4.2。
%
% ^^A ======================================================================
% ^^A Module shared-keylist: 两个类共用的信息字段与常量
% ^^A ======================================================================
%    \begin{macrocode}
%<*shared-keylist>
%    \end{macrocode}
%
% 中文题目的写入器。英文那个只有学位论文类要，留在它自己那半。
%    \begin{macrocode}
\hit@define@title{zh}
%    \end{macrocode}
%
% 两个类都有的信息字段：封面上要填的那批，学院、专业、导师、作者、学号。
%
% \changes{v3.2a}{2026/08/11}{\option{specialty} 改名 \option{major}}
% \changes{v3.2a}{2026/08/11}{加 \option{student-number} 作 \option{student-id} 的别名}
% 专业那个字段 v3.1e 叫 \option{csubject}，注释写的是「中文专业」，v3.2a 起叫
% \option{major}。学科是另一个字段 \option{discipline}（v3.1e 的 \option{cxueke}），
% 封面上「工学博士」的「工学」取的是它，两者不是一回事。
%
% 学号两个名字都对：\option{student-id} 与 \option{student-number}，写哪个都不提示。
%    \begin{macrocode}
\hit@define@term{author-zh}
\hit@define@term{supervisor-zh}
\hit@define@term{affiliation-zh}
\hit@define@term{major-zh}
\hit@define@date@term{date-zh}{zh}{\hit@if@option@TF{bachelor}{day}{month}}
\hit@define@term{student-id}
\hit@define@term@alias{student-number}{student-id}
\hit@declare@info@field
  {
    title-zh , author-zh , supervisor-zh , affiliation-zh , major-zh , date-zh ,
    student-id , student-number
  }
%    \end{macrocode}
%
% 两个类都有的常量。三种存法各一批：\cs{hit@declare@constant} 存进
% \cs{hit@}\meta{名}，\cs{hit@declare@constant@as} 驱动同名的既有命令，
% \cs{hit@declare@constant@ctex} 转发给 \cs{ctexset}。后两批两个类逐字相同。
%    \begin{macrocode}
\hit@declare@constant@as
  {
    equationname , listfigureename , listtableename , listequationename ,
    listequationname , cabstractcname , cabstractename , eabstractcname , eabstractename
  }
\hit@declare@constant@ctex
  {
    chapter-name = chapter/name , appendixname , contentsname , listfigurename ,
    listtablename , figurename , tablename , bibname , indexname
  }
\hit@declare@constant
  {
    algorithm-name-zh , algorithm-name-en , campus-full , authorization-heading ,
    author-signature-field-label , bachelor-affiliation-field-label ,
    bachelor-date-field-label , bachelor-check-date-field-label ,
    bachelor-student-field-label , bachelor-document-type , bachelor-degree-level ,
    bachelor-teacher-comment-field-label , bachelor-teacher-signature-field-label ,
    bibliography-heading-zh , affiliation-field-label-zh ,
    associate-supervisor-field-label-zh , author-field-label-zh ,
    co-supervisor-field-label-zh , date-field-label-zh , degree-of-discipline-format-zh ,
    degree-field-label-zh , school-zh , school-field-label-zh , major-field-label-zh ,
    supervisor-field-label-zh , document-type , thesis-title-prefix , degree-level-zh ,
    date-blank , eqdenote-label-zh , eqdenote-label-en , eqdenote-dash , figure-prefix-en ,
    table-prefix-en , date-field-label-front , graduate-affiliation-field-label ,
    graduate-affiliation-major-field-label , graduate-major-field-label ,
    graduate-date-field-label , graduate-enroll-date-field-label ,
    graduate-student-id-field-label , graduate-student-field-label ,
    graduate-supervisor-field-label , graduate-thesis-field-label ,
    bachelor-school-bottom-mark , harbin-school-bottom-mark , campus , hi ,
    classified-index-international-field-label-zh ,
    classified-index-national-field-label-zh , school-id , school-id-field-label ,
    secret-level-field-label , shenzhen-doctor-interim-note ,
    bachelor-major-field-label , bachelor-student-id-field-label ,
    bachelor-supervisor-field-label , stage-document-type , stage-interim , stage-proposal ,
    supervisor-signature-field-label , title-separator-zh
  }
%    \end{macrocode}
%
% \changes{v3.2a}{2026/08/11}{加 \option{after-}\meta{层级}\option{-number} 与
%   \option{after-}\meta{图表}\option{-caption-number} 两批常量，编号之后的间距
%   从写死改成可配}
% 编号之后排什么。标题四级各一个，图题表题各一个，中英各一份。取值是排版内容，
% 写 \cs{quad}、\cs{enspace} 这种间距最常见，写别的也行。
%
% 默认一律留空，空就退回各处原有的取值，见 \cs{hit@after@number:nnn}。
% \option{after-chapter-number} 只有学位论文类用得上，报告类是 \pkg{ctexart}，没有章。
%    \begin{macrocode}
\hit@declare@constant
  {
    after-chapter-number-zh , after-chapter-number-en ,
    after-section-number-zh , after-section-number-en ,
    after-subsection-number-zh , after-subsection-number-en ,
    after-subsubsection-number-zh , after-subsubsection-number-en ,
    after-figure-caption-number-zh , after-figure-caption-number-en ,
    after-table-caption-number-zh , after-table-caption-number-en
  }
%    \end{macrocode}
%
% 四级标题一起设的写法。写了它等于把四个逐一写一遍，所以后面再单设某一级仍然管用。
%    \begin{macrocode}
\clist_map_inline:nn { zh , en }
  {
    \keys_define:nn { hit / const }
      {
        after-heading-number-#1 .meta:n =
          {
            after-chapter-number-#1       = {##1} ,
            after-section-number-#1       = {##1} ,
            after-subsection-number-#1    = {##1} ,
            after-subsubsection-number-#1 = {##1} ,
          } ,
      }
  }
%    \end{macrocode}
%
% \changes{v3.2a}{2026/08/11}{\option{bib-style} 从 \option{struct} 的部件表挪进
%   \texttt{hit/bib} 族叫 \option{style}。它排不出任何东西，不是部件}
% \changes{v3.2a}{2026/08/16}{\option{style} 再挪进 \option{bibtex} 子组：
%   它是 \file{.bst} 的名字，biblatex 一线用不上（那边锁死国标）}
% \option{bib}/\option{bibtex}/\option{style} 是 \cs{bibliographystyle} 的参数。
% 它原先混在 \option{struct} 的部件表里，可部件说的是「排不排这一块」，样式
% 说的是「排成什么样」，不是一回事。后来又从 \option{bib} 的根层挪进
% \option{bibtex} 子组：挑 \file{.bst} 只在传统一线有意义，biblatex 一线的
% 样式锁死国标（见 \file{hit-bibstyle.dtx}）。键在共享层声明，默认值在这里给。
%
% \changes{v3.2a}{2026/08/11}{\texttt{hit/toc} 族挪到共享层，报告类也能用}
% \texttt{toc} 族。三个键两个类都有，所以声明在共享层：报告一样有目录，一样要
% 排到第几级、西文用不用无衬线体。
%
% 变量在这里声明，早于两个类各自的 \cs{ProcessKeyOptions}，所以旧名的类选项
% 处理器直接写它们就行。取值的含义两个类一致，怎么落到版面上各自决定——
% \option{depth} 数的是级数，学位论文有章、报告没有，所以两边换算到
% \cs{tocdepth} 差一。
%
% \option{english} 只有学位论文类做得了事，报告一份文档只排一份目录，
% 语言跟着类选项 \option{lang} 走。报告类在自己那半把这个键改成一条提示。
%    \begin{macrocode}
\tl_new:N \g__hit_bib_style_tl    \tl_gset:Nn \g__hit_bib_style_tl   { hithesis }
\tl_new:N \g__hit_toc_english_tl  \tl_gset:Nn \g__hit_toc_english_tl { auto }
\int_new:N \g__hit_toc_depth_int  \int_gset:Nn \g__hit_toc_depth_int { 3 }
\bool_new:N \g__hit_toc_latin_sans_bool
\keys_define:nn { hit / toc }
  {
    english .code:n = { \__hit_tristate_set:Nn \g__hit_toc_english_tl {#1} } ,
    depth   .code:n = { \int_gset:Nn \g__hit_toc_depth_int {#1} } ,
  }
\__hit_boolean_key:nnn { toc } { latin-sans } { toc_latin_sans }
%    \end{macrocode}
%
% \changes{v3.2a}{2026/08/11}{\cs{hit@toc@font} 挪到共享层，两个类同一份}
% 目录条目里的西文字体。开关关着时它展开成空，所以可以无条件套在字体那一串
% 前面，不必在每个用到的地方再判断一次。排版时才读标志位，导言区改得动。
%    \begin{macrocode}
\cs_new:Npn \hit@toc@font
  { \bool_if:NT \g__hit_toc_latin_sans_bool { \sffamily } }
%    \end{macrocode}
%
% \changes{v3.2a}{2026/08/13}{加 \texttt{hit/layout} 族}
% \changes{v3.2a}{2026/08/15}{页底策略定名 \option{layout}/\option{bottom}，
%   取值 \texttt{ragged}/\texttt{flush}，与 \hologo{LaTeX} 的 \cs{raggedbottom}
%   与 \cs{flushbottom} 同名。中途用过的 \option{vertical-alignment} 撤销}
% \texttt{layout} 族收 Word 的页面设置。第一个键是 \option{bottom}，取值
% \texttt{ragged}（页面排到哪算哪）与 \texttt{flush}（页底拉齐），名字取自
% \hologo{LaTeX} 自己的两个命令。默认写 \texttt{auto}，眼下等于
% \texttt{ragged}——模板里所有键的默认值都是 \texttt{auto}，这样「没指定」
% 和「指定了某个值」在代码里分得开。
%
% 这一个键同时选定了\emph{按规范的哪一栏排}，因为两者一一对应：
% \begin{center}
% \begin{tabular}{lll}
% \toprule
% & 依据 & 间距 \\\midrule
% \texttt{ragged} & 写作指南 3.2 节的行数与多倍值 & 刚性 \\
% \texttt{flush} & 3.8 节的毫米区间 & 区间即伸缩范围 \\
% \bottomrule
% \end{tabular}
% \end{center}
% 3.2 节给的是确切值，刚性；3.8 节给的全是区间（正文 3\,--\,4\,mm、条标题
% 6\,--\,7\,mm 之类），区间本身就是伸缩范围。按 3.2 节排就没有可拉伸的东西，
% 页底拉不齐；按 3.8 节排就有余地。所以不是两个独立的开关，是一件事的两面。
%
% 反过来说，「按毫米排但保持刚性」这一档不存在——毫米栏从来没给过确切值，
% 硬要刚性就得从区间里挑一个数，那个数没有依据。
%
% 正文的余量算得出来：自然值 3.03\,mm、上限 4\,mm，每行富余 0.97\,mm，
% 一页 33 行合计 32\,mm，而页底缺口通常不到 7\,mm。所以实际每行只拉 0.2\,mm
% 上下，肉眼看不出来。
% \emph{\cs{\_\_hit\_layout\_bottom:n}。}
% \option{layout} 是新版面的接口，不该管旧版面：旧版面进了兼容期就锁死，
% 章节格式在那边设了不管用，页底策略在那边设了却管用，说不通。所以这个键
% 只记下取值，能不能落到标志位上由 \cs{\_\_hit\_layout\_bottom\_apply:} 决定。
%
% 类选项阶段 \cs{g\_\_hit\_new\_layout\_bool} 还没算出来（它要等所有类选项读完
% 才知道 \option{v2-layout} 到底开没开），所以那时按兵不动，等选项后处理
% 那里补一次。用 \cs{hitsetup} 设的时候标志位已经有值，当场就能落。
% \file{hit-report.cls} 没有新旧版面之分，那边这个标志位不存在，直接落。
%
% 旧版面要改页底策略走 \option{v2-layout}\texttt{=\{bottom=\ldots\}}。
%    \begin{macrocode}
\tl_new:N \g__hit_layout_bottom_tl
\cs_new_protected:Npn \__hit_layout_bottom:n #1
  {
    \tl_gset:Nn \g__hit_layout_bottom_tl {#1}
    \__hit_layout_bottom_apply:
  }
\cs_new_protected:Npn \__hit_layout_bottom_apply:
  {
    \tl_if_empty:NF \g__hit_layout_bottom_tl
      { \__hit_layout_bottom_check: }
  }
\cs_new_protected:Npn \__hit_layout_bottom_check:
  {
    \bool_if_exist:NTF \g__hit_layout_two_bool
      {
        \bool_if:NF \g__hit_layout_two_bool
          { \__hit_layout_bottom_set: }
      }
      { \__hit_layout_bottom_set: }
  }
\cs_new_protected:Npn \__hit_layout_bottom_set:
  {
    \str_if_eq:VnTF \g__hit_layout_bottom_tl { ragged }
      { \bool_gset_true:N  \g__hit_raggedbottom_bool }
      { \bool_gset_false:N \g__hit_raggedbottom_bool }
  }
%    \end{macrocode}
%
%
%    \begin{macrocode}
\keys_define:nn { hit / layout }
  {
    bottom .choice: ,
    bottom / auto   .code:n = { \__hit_layout_bottom:n { ragged } } ,
    bottom / ragged .code:n = { \__hit_layout_bottom:n { ragged } } ,
    bottom / flush  .code:n = { \__hit_layout_bottom:n { flush  } } ,
    bottom .value_required:n = true ,
  }
%    \end{macrocode}
%
% \changes{v3.2a}{2026/08/14}{\texttt{layout} 族收下 Word「页面设置」与「文档网格」
%   两个对话框上的键，转给 \cs{hit@msword@layout:n}}
% 其余的键就是 Word 那两个对话框上的格子，一格一个键，换算规则见
% \file{src/utils/hit-msword.dtx}：
% \begin{latex}
% \hitsetup{ layout = {
%   margin-top = 107.75bp , margin-bottom = 85.05bp ,
%   margin-left = 85.05bp , margin-right = 85.05bp ,
%   header = 85.05bp , footer = 65.2bp ,
%   line-pitch = 19.55bp , char-pitch = 12.675bp , size = 12bp ,
% } }
% \end{latex}
%
% 这些键不当场生效，先攒进一张表，一次 \cs{hitsetup} 里的都收齐了再交给
% \cs{hit@msword@layout:n} 重算一遍。版心的每个数都要另外几个数才算得出来——
% 版心高要上下边距，页眉位置要上边距减页眉距边界——一格一格地即算即用会拿
% 半张表去算，还会因为暂时没有行距而误报。
%
% 表是全局只增的，所以分几次写等同于一次写全，后写的盖先写的。
%
% \option{header}/\option{rule} 那两个键与 \option{footer}/\option{shift} 是把类自己的
% 页眉页脚接到 Word 位置上的，来路见 \file{src/utils/hit-msword.dtx}。
%    \begin{macrocode}
\tl_new:N \g__hit_layout_msword_tl
\bool_new:N \l__hit_layout_dirty_bool
\cs_new_protected:Npn \__hit_layout_msword:nn #1#2
  {
    \tl_gput_right:Nx \g__hit_layout_msword_tl { #1 = \exp_not:n {#2} , }
    \bool_set_true:N \l__hit_layout_dirty_bool
  }
%    \end{macrocode}
%
% \changes{v3.2a}{2026/08/16}{\option{layout} 的键按 Word「页面设置」的四个
%   选项卡分组嵌套；\option{head} 撤掉，\option{foot-baseline} 并入
%   \option{footer}/\option{shift}}
% \cs{\_\_hit\_layout\_leaf:nnn}\marg{块}\marg{键}\marg{msword 键} 把一个叶子键
% 定义两遍：一遍在块的子族里，一遍在 \texttt{hit\,/\,layout} 里带全路径。所以
% 下面两种写法等价，嵌套着写和顺手写一条都行：
% \begin{latex}
% layout = { header = { from-edge = 85.05bp } }
% layout = { header/from-edge = 85.05bp }
% \end{latex}
%    \begin{macrocode}
\cs_new_protected:Npn \__hit_layout_leaf:nnn #1#2#3
  {
    \keys_define:nn { hit / layout / #1 }
      {
        #2 .code:n = { \__hit_layout_msword:nn {#3} {##1} } ,
        #2 .value_required:n = true ,
      }
    \keys_define:nn { hit / layout }
      {
        #1 / #2 .code:n = { \__hit_layout_msword:nn {#3} {##1} } ,
        #1 / #2 .value_required:n = true ,
      }
  }
\cs_new_protected:Npn \__hit_layout_block:nn #1#2
  {
    \keys_define:nn { hit / layout / #1 }
      {
        #2 .code:n = { \hit@keys@set:nn { hit / layout / #1 / #2 } {##1} } ,
        #2 .value_required:n = true ,
      }
    \keys_define:nn { hit / layout }
      {
        #1 / #2 .code:n = { \hit@keys@set:nn { hit / layout / #1 / #2 } {##1} } ,
        #1 / #2 .value_required:n = true ,
      }
  }
\clist_map_inline:nn { paper , margin , header , footer , grid }
  {
    \keys_define:nn { hit / layout }
      {
        #1 .code:n = { \hit@keys@set:nn { hit / layout / #1 } {##1} } ,
        #1 .value_required:n = true ,
      }
  }
\__hit_layout_leaf:nnn { paper } { width }  { paper-width }
\__hit_layout_leaf:nnn { paper } { height } { paper-height }
\__hit_layout_leaf:nnn { margin } { top }    { margin-top }
\__hit_layout_leaf:nnn { margin } { bottom } { margin-bottom }
\__hit_layout_leaf:nnn { margin } { left }   { margin-left }
\__hit_layout_leaf:nnn { margin } { right }  { margin-right }
\__hit_layout_leaf:nnn { header } { from-edge } { header }
\__hit_layout_leaf:nnn { footer } { from-edge } { footer }
\__hit_layout_block:nn { header } { rule }
\__hit_layout_block:nn { footer } { shift }
\__hit_layout_block:nn { grid } { chars }
\__hit_layout_block:nn { grid } { lines }
\__hit_layout_leaf:nnn { header / rule } { width }     { header-rule-width }
\__hit_layout_leaf:nnn { header / rule } { from-text } { header-rule-from-text }
\__hit_layout_leaf:nnn { footer / shift } { vertical }   { footer-shift-vertical }
\__hit_layout_leaf:nnn { footer / shift } { horizontal } { footer-shift-horizontal }
\__hit_layout_leaf:nnn { grid / chars } { pitch }    { char-pitch }
\__hit_layout_leaf:nnn { grid / chars } { per-line } { chars-per-line }
\__hit_layout_leaf:nnn { grid / lines } { pitch }    { line-pitch }
\__hit_layout_leaf:nnn { grid / lines } { per-page } { lines-per-page }
%    \end{macrocode}
%
% \changes{v3.2a}{2026/08/15}{加 \option{layout}/\option{char-pitch-en}，
%   字距增量也加到西文字母上}
% \changes{v3.2a}{2026/08/16}{\option{char-pitch-en} 改名
%   \option{grid}/\option{chars}/\option{latin}}
% \option{grid}/\option{chars}/\option{latin} 决定网格字距加不加到西文字母上。
% Word 那边 \texttt{charSpace} 是加在\emph{每一个}字符后面的，汉字、字母、数字、
% 空格一视同仁——实测每个西文字母后面也多 0.4543\,bp，一行 82 字累计 37\,bp，
% 差不多三个汉字宽。默认关，因为它会动到参考文献、URL 与代码环境里的每一个
% 西文串。
%
% 原先叫 \option{char-pitch-en}，那个名字像是有个 \option{char-pitch-zh} 与它配对，
% 其实没有：汉字那一侧不是全文开关，是 \option{format} 里的
% \option{snap-to-grid}/\option{horizontal}，按元素发 \cs{CJKglue}。两边不对称是
% 实现逼出来的——\cs{CJKglue} 是个宏，能按组换；西文那边要发
% \cs{defaultfontfeatures} 的 \texttt{LetterSpace}，得重新声明整个字族，
% 只能全文来一次。所以它留在 \option{layout}。
%    \begin{macrocode}
\bool_new:N \g__hit_layout_pitch_en_bool
\__hit_boolean_key:nnn { layout / grid / chars } { latin } { layout_pitch_en }
\__hit_boolean_key:nnn { layout } { grid / chars / latin } { layout_pitch_en }
%    \end{macrocode}
%
%    \begin{macrocode}
\cs_new_protected:Npn \hit@layout@set:n #1
  {
    \bool_set_false:N \l__hit_layout_dirty_bool
    \hit@keys@set:nn { hit / layout } {#1}
    \bool_if:NT \l__hit_layout_dirty_bool
      { \exp_args:NV \hit@msword@layout:n \g__hit_layout_msword_tl }
  }
%    \end{macrocode}
%
% \changes{v3.2a}{2026/08/12}{\pkg{etoolbox} 从 thesis-keylist 挪进 shared-keylist，
%   两个类共用一次加载。原先两个参考文献模块各装一次，那两次一个宏都没用上}
% \pkg{etoolbox} 在这里装。下面 \cs{AfterPreamble} 与 \cs{ifdefempty} 要它，
% \file{src/utils/hit-caption.dtx} 的三处 \cs{patchcmd} 也要。两个消费者都排在
% \cs{LoadClass} 之前，\file{src/flow/hit-*-deps.dtx} 太晚，所以在这儿。
% \pkg{etoolbox} 不牵扯字体与基类，提前装没有副作用。
%    \begin{macrocode}
\RequirePackage{etoolbox}
%</shared-keylist>
%    \end{macrocode}
%
%
% \section{学位论文类的键清单（keylist）}
%
% \changes{v3.2a}{2026/08/10}{元信息字段的声明单独成文件。它说的是「有哪些字段」，
%   与 \file{src/setup/} 说的「字段取什么值」是一对，都是数据；
%   字段怎么存、旧名怎么兼容那套机制在 \file{src/utils/hit-meta.dtx}}
% 学位论文用到的元信息字段：字段名、旧名映射、\cs{hitsetup} 各族的路由，
% 以及摘要环境。都是清单，不含机制。
%
% ^^A ======================================================================
% ^^A Module thesis-keylist: 元信息字段（\hitsetup 接受的 key）（hit-thesis）
% ^^A ======================================================================
%    \begin{macrocode}
%<*thesis-keylist>
% \changes{v3.2a}{2026/08/10}{\cs{hitsetup} 的底层机制挪到 \file{src/utils/hit-meta.dtx}}

\hit@define@term{secret-level} %密级
\hit@define@term{classified-index-national}  %国内图书分类号
\hit@define@term{classified-index-international}  %国际图书分类号

\hit@define@term{discipline-zh} %中文学科
\hit@define@term{associate-supervisor-zh} %中文副导师
\hit@define@term{co-supervisor-zh}%中文联合导师
%    \end{macrocode}
%
% \changes{v3.1b}{2022/06/06}{增加深圳封面最下方的时间信息}
% \changes{v3.1d}{2025/03/03}{根据 \issue[226]，将 \texttt{szshortcdate} 改为 \texttt{firstpagecdate}，只要赋值就会显示}
% 根据 \issue[95] 增加一个变量 \texttt{szshortcdate} 来存储答辩时间。
% 根据 \issue[226]，将 \texttt{szshortcdate} 改为 \texttt{firstpagecdate}，只要赋值就会显示。
%    \begin{macrocode}
\hit@define@date@term{cover-date-zh}{zh}{\hit@if@option@TF{bachelor}{day}{month}}
\hit@define@term{report-number}
\hit@define@term{station-zh}
%    \end{macrocode}
%
% 添加博后封皮选项
% \changes{v3.0.12}{2020/11/02}{添加博后封皮选项}
%    \begin{macrocode}
\hit@define@date@term{work-finish-date-zh}{zh}{auto}
\hit@define@date@term{report-submit-date-zh}{zh}{auto}
\hit@define@date@term{research-start-date-zh}{zh}{auto}
\hit@define@date@term{research-end-date-zh}{zh}{auto}
%    \end{macrocode}
%
% 添加博后封皮的title选项
% \changes{v3.0.0}{2020/05/21}{添加博后封皮的title选项}
%    \begin{macrocode}

\hit@define@term{student-type-zh}%



\hit@define@title{en}
\hit@define@term{discipline-en} %英文学科
\hit@define@term{author-en} %英文作者
\hit@define@term{supervisor-en} %英文导师
\hit@define@term{associate-supervisor-en} %英文副导师
\hit@define@term{co-supervisor-en} %英文联合导师
\hit@define@term{affiliation-en}
\hit@define@term{major-en}
\hit@define@date@term{date-en}{en}{month}
\hit@define@term{student-type-en}
%    \end{macrocode}
%
% \changes{v3.2a}{2026/08/12}{四个摘要环境从 \pkg{environ} 的 \cs{Collect@Body} 改用
%   内核的 \texttt{+b} 参数类型。\pkg{environ} 至此没人用，已从两个类里删掉}
% 摘要是「先收着，后面再排」：用户在正文位置写 \env{abstract-zh}，内容存进
% \cs{hit@abstract@zh}，真正排版在摘要页那边。所以要把整个环境体当参数抓下来。
% \cs{NewDocumentEnvironment} 的 \texttt{b} 参数类型就是环境体，必须排在最后，
% 前面的 \texttt{+} 表示允许 \cs{par}。
%
% 环境体是当 token 抓的，里面不能放 \cs{verb} 或 \env{lstlisting}。
% 原先的 \cs{Collect@Body} 是一样的限制。
%    \begin{macrocode}
\cs_new_protected:Npn \hit@collect@abstract@zh #1 { \cs_gset:Npn \hit@abstract@zh {#1} }
\cs_new_protected:Npn \hit@collect@abstract@en #1 { \cs_gset:Npn \hit@abstract@en {#1} }
\NewDocumentEnvironment { abstract-zh } { +b } { \hit@collect@abstract@zh {#1} } { }
\NewDocumentEnvironment { abstract-en } { +b } { \hit@collect@abstract@en {#1} } { }
\NewDocumentEnvironment { cabstract } { +b }
  {
    \hit@warn@renamed@environment { cabstract } { abstract-zh }
    \hit@collect@abstract@zh {#1}
  }
  { }
\NewDocumentEnvironment { eabstract } { +b }
  {
    \hit@warn@renamed@environment { eabstract } { abstract-en }
    \hit@collect@abstract@en {#1}
  }
  { }
%    \end{macrocode}
%
% \changes{v3.2a}{2026/08/10}{关键词的逗号拆分挪到 \file{src/utils/hit-meta.dtx}}
% 关键词的拆分机制见 \file{src/utils/hit-meta.dtx}。
%    \begin{macrocode}
\hit@parse@keywords{keywords-zh}{keywords@separator@zh}
\hit@parse@keywords{keywords-en}{keywords@separator@en}
%    \end{macrocode}
%
% \changes{v3.2a}{2026/08/10}{\cs{hitsetup} 改用 l3keys 做入口，加一层按职责
%   划分的子族路由。子族只做转发，键由 \texttt{hit/class} 族处理，行为
%   与改动前一致；后续再逐族把键迁过来}
% \cs{hitsetup} 的入口换成 l3keys。根族只做路由，把 \option{info}、\option{style}
% 等分派到对应子族；子族一律把键原样转回 \texttt{hit/class}
% 族，所以两种写法结果相同：
% \begin{latex}
% \hitsetup{title-zh={...}}              % 扁平写法，一直支持
% \hitsetup{info={title-zh={...}}}       % 新的分族写法
% \end{latex}
% 根族的 \option{unknown} 兜底负责扁平写法，子族的兜底负责分族写法。等某个键真正
% 迁到 l3keys 之后，它就不再走兜底，两条路自然收敛。兜底本身、三态取值与布尔键的
% 声明机制都在 \file{src/utils/hit-meta.dtx}。
%
% \changes{v3.2a}{2026/08/10}{元信息字段正式注册进 \texttt{hit/info} 族，
%   不再依赖子族兜底转发}
% 元信息字段名单。这批由 \cs{hit@define@term} 与 \cs{hit@parse@keywords} 批量生成，
% 每个字段同时有一个同名命令（\cs{title-zh} 等），键的处理器直接调用那个命令，
% 所以三种写法结果一致：\cs{hitsetup} 的扁平写法、分族写法、以及命令写法。
%    \begin{macrocode}
\hit@declare@info@field
  {
    secret-level , classified-index-national , classified-index-international ,
    discipline-zh , associate-supervisor-zh , co-supervisor-zh , cover-date-zh ,
    report-number , station-zh , work-finish-date-zh , report-submit-date-zh ,
    research-start-date-zh , research-end-date-zh , student-type-zh , title-en ,
    discipline-en , author-en , supervisor-en , associate-supervisor-en ,
    co-supervisor-en , affiliation-en , major-en , date-en , student-type-en ,
    keywords-zh , keywords-en
  }
%    \end{macrocode}
%
% \changes{v3.2a}{2026/08/10}{\file{cfg} 里的可本地化字符串做成 \texttt{hit/const}
%   族的键，用户在导言区即可改写，不必去动生成物}
% 可本地化字符串的名单。默认值在类文件末尾给出，\cs{hitsetup} 在导言区执行，
% 所以用户设的值覆盖默认值。键名怎么由宏名得来见 \file{src/utils/hit-meta.dtx}。
%    \begin{macrocode}
\hit@declare@constant
  {
    bachelor-copyright-heading , bachelor-authorization-heading ,
    bachelor-declarations-heading-line-one , bachelor-declarations-heading-line-two ,
    bachelor-declarations-heading-short , acknowledgement-heading-zh ,
    acknowledgement-heading-en , declarations-heading-zh , declarations-heading-en ,
    authorization-text , bachelor-copyright-text , bibliography-heading-en ,
    brace-left-zh , brace-right-zh , keywords-separator-zh , keywords-field-label-zh ,
    conclusion-heading-zh , keywords-field-label-en , contents-page-heading ,
    contents-table-heading , contents-figure-heading , nomenclature-heading , conclusion-heading-en ,
    corresponding-address-heading-zh , corresponding-address-heading-en ,
    copyright-heading , bachelor-declarations-heading , copyright-text ,
    student-type-academic-zh , student-type-academic-en ,
    student-type-professional-zh , student-type-professional-en ,
    denotation-heading-zh , denotation-heading-en , doctor-publication-heading-zh ,
    doctor-publication-heading-en , affiliation-field-label-en ,
    associate-supervisor-field-label-en , author-field-label-en , brace-left-en ,
    brace-right-en , co-supervisor-field-label-en , date-field-label-en ,
    degree-of-discipline-format-en , degree-field-label-en , keywords-separator-en ,
    school-en , school-field-label-en , major-field-label-en ,
    supervisor-field-label-en , degree-level-en , degree-level-attribute , heading-spread ,
    index-heading-en , index-heading-zh , classified-index-international-field-label-en ,
    classified-index-national-field-label-en , postdoc-udc-field-label ,
    postdoc-author-field-label , postdoc-classified-index-field-label , postdoc-report ,
    postdoc-research-end-date-field-label , postdoc-work-finish-date-field-label ,
    postdoc-major-field-label , postdoc-number-field-label ,
    postdoc-station-field-label , postdoc-secret-level-field-label , postdoc-separator ,
    postdoc-research-start-date-field-label , postdoc-report-submit-date-field-label ,
    postdoc-co-supervisor-field-label , publication-heading-zh , publication-heading-en ,
    defense-heading-zh , defense-heading-en , resume-heading-zh , resume-heading-en ,
    defense-reviewer-label , defense-name-label , defense-title-label ,
    defense-affiliation-label , defense-discipline-label , defense-role-label ,
    defense-committee-label , defense-chair-label , defense-member-label ,
    defense-secretary-label , defense-resolution-label ,
    bachelor-authorization-text , bachelor-thesis-field-label , title-separator-en
  }
%    \end{macrocode}
%
% \changes{v3.2a}{2026/08/10}{参考文献的五个选项迁入 \texttt{hit/bib} 族，
%   改用连字符命名。旧名转回 \texttt{hit/class} 族处理，行为不变}
% \changes{v3.2a}{2026/08/10}{排版开关迁入 \texttt{hit/style} 族并改用连字符命名}
% 两个族的键。排版开关都是普通布尔，用 \cs{\_\_hit\_style\_pair:nn} 成对声明，
% 新名设的是旧标志位本身，所以两套名字天然一致。
% \option{heading-latin-sans} 跟共享层的 \option{toc}/\option{latin-sans} 是一对，
% 一个管章节标题里的西文一个管目录里的。参考文献的五个键此前分散在
% \file{book} 与 \file{art} 两个模块、命名也不统一（\texttt{citetwo} 看不出是
% 控制什么的），新旧两套名字设的是同一批内部变量，所以可以混用。
%    \begin{macrocode}
\keys_define:nn { hit / style }
  {
    ziju .choices:nn = { auto , regu , word }
      {
        \str_if_eq:VnF \l_keys_choice_tl { auto }
          {
            \bool_gset_false:N \g__hit_ziju_regu_bool
            \bool_gset_false:N \g__hit_ziju_word_bool
            \bool_gset_true:c { g__hit_ziju_ \l_keys_choice_tl _bool }
          }
      } ,
    ziju .value_required:n = true ,
    header-rule .code:n = { \__hit_tristate_set:Nn \g__hit_header_rule_tl {#1} } ,
    page-number-line .code:n =
      { \hit@keys@set:nn { hit / style / page-number-line } {#1} } ,
  }
\keys_define:nn { hit / style / page-number-line }
  {
    mainmatter  .code:n = { \__hit_tristate_set:Nn \g__hit_page_number_line_main_tl  {#1} } ,
    frontmatter .code:n = { \__hit_tristate_set:Nn \g__hit_page_number_line_front_tl {#1} } ,
  }
\clist_map_inline:nn
  {
    { title-en-auto } { etitleauto } ,
    { open-right } { openright } ,
    { library } { library } ,
  }
  { \__hit_style_pair:nn #1 }
%    \end{macrocode}
%
% \changes{v3.2a}{2026/08/10}{加 \texttt{hit/math} 族与
%   \option{math/numbering-parenthesis-style}}
% \texttt{math} 族。\option{numbering-parenthesis-style} 取 \texttt{full} 是全角
% 括号、\texttt{half} 是半角，默认全角。英文论文不受它影响，一律半角，那个判断在
% \cs{hit@equation@paren@left} 里。
%    \begin{macrocode}
\keys_define:nn { hit / math }
  {
    numbering-parenthesis-style .choice: ,
    numbering-parenthesis-style / full .code:n =
      { \bool_gset_true:N \g__hit_equation_paren_full_bool } ,
    numbering-parenthesis-style / half .code:n =
      { \bool_gset_false:N \g__hit_equation_paren_full_bool } ,
    numbering-parenthesis-style / auto .code:n = { } ,
    numbering-parenthesis-style .value_required:n = true ,
  }
%    \end{macrocode}
%
% \changes{v3.2a}{2026/08/16}{参考文献的键按 kebab-biblatex-gb7714 的词汇重起：
%   \option{compress-min}、\option{uppercase-family}、\option{max-bib-names}
%   （另扩一个 \option{max-cite-names}）、\option{numbering-align}、
%   \option{split-entry}。\option{style} 挪进
%   \option{bibtex} 子组（见 \file{hit-bibstyle.dtx}），中途用过的
%   \option{adjacent-cite}、\option{author-case}、\option{maximum-author}、
%   \option{label-align}、\option{split-item} 没进过发布版，直接撤}
% \option{compress-min} 是「相邻连续序号几个起压成区间」：\texttt{2} 排
% \texttt{[1-2]}（默认），\texttt{3} 排 \texttt{[1,2]}，三个起总归压。两条线
% 都认，落点都是官方机制：传统一线转发给 \pkg{gbt7714.sty} 的
% \texttt{mincompress}（v3.2a 起，原先自家打的 \pkg{natbib} 补丁退役），
% biblatex 一线落到 \pkg{gb7714-2015} 的 \texttt{gbrefcompress} 计数器。
% 两边都不锁时机，随设随生效；旧类选项 \option{citetwo} 的那对标志位并进
% 同一个取值函数，两条线都认它。
%
% \option{uppercase-family}（\texttt{true} 著者姓名全大写即国标排法、
% \texttt{false} 保持 \file{.bib} 原样）、\option{max-bib-names}（文献表里
% 最多列几位作者，\texttt{0} 不限）、\option{max-cite-names}（正文引用侧的
% 同款，数字引用不出姓名，只在 biblatex 一线对 \cs{textcite} 这类起作用）、
% \option{numbering-align}（标号 \texttt{left}/\texttt{right}）这几个在
% biblatex 一线是装载时换算成它的选项的，装载之后再设不生效，设值器前面搭
% 一句守卫（见 \file{hit-bibstyle.dtx}）。\option{split-entry} 排版时才读，
% 不用守。
%    \begin{macrocode}
\int_new:N \g__hit_cite_compress_int
\cs_new:Npn \__hit_bib_compress_effective:
  {
    \int_compare:nNnTF \g__hit_cite_compress_int > 0
      { \int_use:N \g__hit_cite_compress_int }
      {
        \bool_if:NTF \g__hit_citetwo_comma_bool { 3 }
          { \bool_if:NTF \g__hit_citetwo_endash_bool { 2 } { 0 } }
      }
  }
\cs_new_protected:Npn \hit@bib@cite@compress@apply:
  {
    \int_compare:nNnF { \__hit_bib_compress_effective: } = 0
      {
        \@ifpackageloaded { biblatex }
          {
            \cs_if_exist:cT { c@gbrefcompress }
              { \setcounter { gbrefcompress } { \__hit_bib_compress_effective: } }
          }
          {
            \@ifpackageloaded { gbt7714 }
              {
                \__hit_bib_if_mincompress:T
                  {
                    \use:e
                      {
                        \exp_not:N \keys_set:nn { gbt7714 }
                          { mincompress = \__hit_bib_compress_effective: }
                      }
                  }
              } { }
          }
      }
  }
%    \end{macrocode}
%
% \changes{v3.2a}{2026/08/16}{转发 \texttt{mincompress} 前先查这个键在不在，
%   \pkg{gbt7714}~v2 没有它}
% \texttt{mincompress} 是 \pkg{gbt7714}~v3 才有的键。\TeX~Live 2025 装的是
% v2.1.5（2022/10/03），\TeX~Live 2026 才换成 v3.0.1（2026/07/20）。直接
% \cs{keys\_set:nn} 过去，v2 上会报 \texttt{The key 'gbt7714/mincompress' is
% unknown}，配上 \texttt{-halt-on-error} 就是编不出来。
%
% 判据用 \cs{keys\_if\_exist:nnTF} 查键本身，不查包的版本号或日期：版本号将来
% 怎么走不好说，键在不在才是我们真正依赖的东西。这个判断在 \TeX~Live 2025 与
% 2026 的 \pkg{l3kernel} 里都有；万一更老的树上没有，外面那层
% \cs{cs\_if\_exist:NTF} 会让它退回「不转发」，仍然编得出来，只是压缩阈值用
% \pkg{gbt7714} 自己的默认值。
%    \begin{macrocode}
\prg_new_protected_conditional:Npnn \__hit_bib_if_mincompress: { T }
  {
    \cs_if_exist:NTF \keys_if_exist:nnTF
      {
        \keys_if_exist:nnTF { gbt7714 } { mincompress }
          { \prg_return_true: } { \prg_return_false: }
      }
      { \prg_return_false: }
  }
%    \end{macrocode}
%
%    \begin{macrocode}
\keys_define:nn { hit / bib }
  {
    compress-min .choices:nn = { auto , 2 , 3 }
      {
        \str_case:Vn \l_keys_choice_tl
          {
            { auto } { \int_gset:Nn \g__hit_cite_compress_int { 0 } }
            { 2 }
              {
                \int_gset:Nn \g__hit_cite_compress_int { 2 }
                \bool_gset_false:N \g__hit_citetwo_comma_bool
                \bool_gset_true:N  \g__hit_citetwo_endash_bool
              }
            { 3 }
              {
                \int_gset:Nn \g__hit_cite_compress_int { 3 }
                \bool_gset_true:N  \g__hit_citetwo_comma_bool
                \bool_gset_false:N \g__hit_citetwo_endash_bool
              }
          }
        \hit@bib@cite@compress@apply:
      } ,
    compress-min .value_required:n = true ,
    numbering-align .code:n =
      {
        \str_if_eq:nnF {#1} { auto }
          {
            \hit@bib@key@guard:n { numbering-align }
            \str_gset:Nn \g_bib_labelalign_str {#1}
          }
      } ,
    uppercase-family .choices:nn = { true , false , auto }
      {
        \str_if_eq:VnF \l_keys_choice_tl { auto }
          {
            \hit@bib@key@guard:n { uppercase-family }
            \str_if_eq:VnTF \l_keys_choice_tl { true }
              { \bool_gset_true:N  \g_bib_authorupper_bool }
              { \bool_gset_false:N \g_bib_authorupper_bool }
          }
      } ,
    uppercase-family .default:n = true ,
    max-bib-names .code:n =
      {
        \str_if_eq:nnF {#1} { auto }
          {
            \hit@bib@key@guard:n { max-bib-names }
            \int_gset:Nn \g_bib_maxauthor_int {#1}
            \int_if_zero:nT { \g_bib_maxauthor_int }
              { \int_gdecr:N \g_bib_maxauthor_int }
          }
      } ,
    max-cite-names .code:n =
      {
        \str_if_eq:nnF {#1} { auto }
          {
            \hit@bib@key@guard:n { max-cite-names }
            \int_gset:Nn \g__hit_bib_maxcite_int {#1}
            \int_compare:nNnT \g__hit_bib_maxcite_int = 0
              { \int_gdecr:N \g__hit_bib_maxcite_int }
          }
      } ,
    split-entry .choice: ,
    split-entry / true  .code:n = { \bool_gset_true:N \g__hit_splitbibitem_bool } ,
    split-entry / false .code:n = { \bool_gset_false:N \g__hit_splitbibitem_bool } ,
    split-entry / auto  .code:n = { } ,
    split-entry .default:n = true ,
  }
%    \end{macrocode}
%
% \changes{v3.2a}{2026/08/10}{加旧名到新名的映射表与废弃警告。旧名照常生效，
%   只是编译时提示新写法}
% 兼容期的处理。旧名一律照常生效，只在日志里给一条提示。三类提示分开：
% 改过名的选项逐个报，元信息字段的扁平写法合并成一条报，
% \cs{documentclass} 里传排版类选项单独报。
%
% \changes{v3.2a}{2026/08/11}{\option{tocblank} 废弃：仍然认这个键，但不再起作用}
% 另有一张废弃表，记的是没有新名可改的选项。这类键照旧认，只是什么也不做，
% 提示里说清为什么。\option{tocblank} 插的那个空行规范里没有，所以整条去掉，
% 行为等同于原先的 \option{tocblank}\texttt{=false}。
% \changes{v3.2a}{2026/08/15}{\option{glue} 与 \option{raggedbottom} 从改名表
%   挪到拆分表：这两个不是换了名字，是按新旧版面分成了两个去处}
% \option{glue} 和 \option{raggedbottom} 不能只报一个新名字。它们原来同时
% 承担两件事：调旧版面那套手写的伸缩量，以及决定页底排不排齐。新版面把这
% 两件事分开了，所以得问清楚用户到底要哪一种，两条路都写给他看。
%    \begin{macrocode}
\prop_gset_from_keyval:Nn \g__hit_split_option_prop
  {
    glue         = { split } ,
    raggedbottom = { split } ,
  }
%    \end{macrocode}
%
%    \begin{macrocode}
\prop_gset_from_keyval:Nn \g__hit_removed_option_prop
  {
    tocblank = { 目录里第一章之前不再插空行，规范里没有这一条 } ,
  }
\prop_gset_from_keyval:Nn \g__hit_renamed_option_prop
  {
    arialtoc         = { toc = { latin-sans = ... } } ,
    tocfour          = { toc = { depth = 4 } } ,
    engtoc           = { toc = { english = ... } } ,
    arialtitle       = { format = { heading = { font-en = sans } } } ,
    chapterbold      = { format = { chapter = { bold = ... } } } ,
    chapterhang      = { format = { chapter = { indent-special = hanging } } } ,
    capcenterlast    = { format = { caption = { alignment = last-line-centered } } } ,
    subcapcenterlast = { format = { caption-sub = { alignment = last-line-centered } } } ,
    absupper         = { format = { abstract-heading = { all-caps = ... } } } ,
    etitleauto       = { style = { title-en-auto = ... } } ,
    openright        = { style = { open-right = ... } } ,
    library          = { style = { library = ... } } ,
    zijv             = { style = { ziju = ... } } ,
    bsheadrule       = { style = { header-rule = ... } } ,
    bsmainpagenumberline  = { style = { page-number-line = { mainmatter = ... } } } ,
    bsfrontpagenumberline = { style = { page-number-line = { frontmatter = ... } } } ,
    citetwo          = { bib = { compress-min = 2~|~3 } } ,
    splitbibitem     = { bib = { split-entry = ... } } ,
    biblabelalign    = { bib = { numbering-align = ... } } ,
    bibauthorupper   = { bib = { uppercase-family = true } } ,
    bibmaxauthor     = { bib = { max-bib-names = ... } } ,
  }
%    \end{macrocode}
%
% \changes{v3.2a}{2026/08/10}{元信息字段的旧名走废弃期：仍然可用，用一次提示一次}
% 旧名到新名的对照。旧名在 v3.1e 就是公开接口，写在用户的 \file{cover.tex} 里，
% 所以保留到 v4。
%    \begin{macrocode}
\hit@define@legacy@terms
  {
    statesecrets=secret-level ,
    natclassifiedindex=classified-index-national ,
    intclassifiedindex=classified-index-international ,
    ctitle=title-zh ,
    cxueke=discipline-zh ,
    cauthor=author-zh ,
    csupervisor=supervisor-zh ,
    cassosupervisor=associate-supervisor-zh ,
    ccosupervisor=co-supervisor-zh ,
    caffil=affiliation-zh ,
    csubject=major-zh ,
    firstpagecdate=cover-date-zh ,
    cdate=date-zh ,
    cnumber=report-number ,
    cpositionname=station-zh ,
    cfinishdate=work-finish-date-zh ,
    csubmitdate=report-submit-date-zh ,
    cstartdate=research-start-date-zh ,
    cenddate=research-end-date-zh ,
    cstudentid=student-id ,
    cstudenttype=student-type-zh ,
    etitle=title-en ,
    exueke=discipline-en ,
    eauthor=author-en ,
    esupervisor=supervisor-en ,
    eassosupervisor=associate-supervisor-en ,
    ecosupervisor=co-supervisor-en ,
    eaffil=affiliation-en ,
    esubject=major-en ,
    edate=date-en ,
    estudenttype=student-type-en ,
    ckeywords=keywords-zh ,
    ekeywords=keywords-en
  }
\hit@define@legacy@title{ctitlecover}{zh}{first}
\hit@define@legacy@title{ctitleone}{zh}{first}
\hit@define@legacy@title{ctitletwo}{zh}{second}
\hit@define@legacy@title{csubtitle}{zh}{subtitle}
\hit@define@legacy@title{esubtitle}{en}{subtitle}
\prop_gset_from_keyval:Nn \g__hit_renamed_load_option_prop
  {
    font        = { font-en = ... } ,
    cjk-font    = { font-zh = ... } ,
    type        = { degree-level } ,
    newgeometry = { v2-layout~=~{~bachelor~,~graduate~=~one|two~} } ,
  }
%    \end{macrocode}
%
% \changes{v3.2a}{2026/08/10}{提示文案、只报一次的机制与类选项扫描挪到
%   \file{src/utils/hit-meta.dtx}}
% 单双面由学位与校区决定，所以扫描类选项时额外查一下 \option{oneside} 与
% \option{twoside}。另外两个类没有这条。
%    \begin{macrocode}
\cs_set_protected:Npn \__hit_class_option_extra:
  {
    \clist_if_in:nVT { oneside , twoside } \l__hit_warn_message_tl
      {
        \exp_args:NV \__hit_warn_once:nnn \l__hit_warn_message_tl
          { manual-sides } { }
      }
  }
%    \end{macrocode}
%
% \changes{v3.2a}{2026/08/12}{加 \texttt{defense} 族：评阅人与答辩委员会名单、
%   答辩决议正文。原先答辩决议那一页只能贴扫描件，现在可以直接填}
% \texttt{defense} 族收答辩决议那一页要的数据。一个人一条记录，字段具名：
% \begin{latex}
% \hitsetup{ defense = {
%   reviewer = {name=张三, title=教授（博导）, affiliation=哈尔滨工业大学,
%               discipline=控制科学与工程},
%   reviewer = {name=李四, …},
%   chair     = {…},
%   member    = {…},
%   secretary = {…},
%   resolution = {答辩委员会经过讨论，一致同意通过答辩，建议授予工学博士学位。},
% } }
% \end{latex}
%
% \textbf{同一个键写多遍是累加，不是覆盖。} \option{reviewer} 与 \option{member}
% 靠这个收任意条。\pkg{l3keys} 按书写顺序逐条执行，键的代码是「往队列右边追加」，
% 写几遍就是几个人，顺序也就是写的顺序。
%
% 四个名单各一个队列，与纸质表的分块一一对应。每一项存的是那个人的整串键值，
% 排版时再解开，所以少写一个字段只是那一格空着，不会像并列列表那样把后面全顶错位。
%    \begin{macrocode}
\seq_new:N \g__hit_defense_reviewer_seq
\seq_new:N \g__hit_defense_chair_seq
\seq_new:N \g__hit_defense_member_seq
\seq_new:N \g__hit_defense_secretary_seq
\tl_new:N  \g__hit_defense_resolution_tl
\keys_define:nn { hit / defense }
  {
    reviewer   .code:n = { \seq_gput_right:Nn \g__hit_defense_reviewer_seq  {#1} } ,
    chair      .code:n = { \seq_gput_right:Nn \g__hit_defense_chair_seq     {#1} } ,
    member     .code:n = { \seq_gput_right:Nn \g__hit_defense_member_seq    {#1} } ,
    secretary  .code:n = { \seq_gput_right:Nn \g__hit_defense_secretary_seq {#1} } ,
    resolution .code:n = { \tl_gset:Nn \g__hit_defense_resolution_tl {#1} } ,
    unknown .code:n =
      { \msg_warning:nnV { hithesis } { unknown-key } \l_keys_key_str } ,
  }
%    \end{macrocode}
%
% 一条记录里的四个字段。排版时把一项喂进这个族，四个变量就有了值；
% 每次解开之前先清空，免得上一个人的值漏到下一个人身上。
%    \begin{macrocode}
\tl_new:N \l__hit_defense_name_tl
\tl_new:N \l__hit_defense_title_tl
\tl_new:N \l__hit_defense_affiliation_tl
\tl_new:N \l__hit_defense_discipline_tl
\keys_define:nn { hit / defense / person }
  {
    name        .tl_set:N = \l__hit_defense_name_tl ,
    title       .tl_set:N = \l__hit_defense_title_tl ,
    affiliation .tl_set:N = \l__hit_defense_affiliation_tl ,
    discipline  .tl_set:N = \l__hit_defense_discipline_tl ,
    unknown .code:n =
      { \msg_warning:nnV { hithesis } { unknown-key } \l_keys_key_str } ,
  }
\cs_new_protected:Npn \hit@defense@unpack:n #1
  {
    \tl_clear:N \l__hit_defense_name_tl
    \tl_clear:N \l__hit_defense_title_tl
    \tl_clear:N \l__hit_defense_affiliation_tl
    \tl_clear:N \l__hit_defense_discipline_tl
    \hit@keys@set:nn { hit / defense / person } {#1}
  }
%    \end{macrocode}
%    \begin{macrocode}
\__hit_check_class_options:
\keys_define:nn { hit }
  {
    info  .code:n = { \hit@keys@set:nn { hit / info } {#1} } ,
    style .code:n = { \hit@keys@set:nn { hit / style } {#1} } ,
    layout .code:n = { \hit@layout@set:n {#1} } ,
    format .code:n = { \hit@keys@set:nn { hit / format } {#1} } ,
    toc   .code:n = { \hit@keys@set:nn { hit / toc } {#1} } ,
    bib   .code:n =
      { \hit@keys@set:nn { hit / bib } {#1} } ,
    const .code:n = { \hit@keys@set:nn { hit / const } {#1} } ,
    math  .code:n = { \hit@keys@set:nn { hit / math } {#1} } ,
    struct .code:n =
      { \__hit_structure_reset_written: \hit@keys@set:nn { hit / struct } {#1} } ,
    defense .code:n = { \hit@keys@set:nn { hit / defense } {#1} } ,
    thm .code:n = { \hit@theorem@style {#1} } ,
    unknown .code:n = { \__hit_legacy_key:n {#1} } ,
  }
\RenewDocumentCommand \hitsetup { O{} +m }
  {
    \hit@if@condition:nT {#1}
      { \hit@keys@set:nn { hit } {#2} \hit@bib@backend@apply: }
  }
%    \end{macrocode}
%
% \changes{v3.1d}{2025/03/03}{将 \cs{firstpagecdate} 的设置藏起来}
% \changes{v3.2a}{2026/08/13}{封面日期的回退改看存原值那个宏。日期字段改成
%   格式化取值之后，\cs{hit@cover@date@zh} 恒非空，\cs{ifdefempty} 那条判断
%   永远走不到回退，首页日期整个空掉}
% 首页那行日期没单独写就跟着封面日期走。判断看的是存原值的
% \cs{hit@cover@date@zh@raw}——\cs{hit@cover@date@zh} 现在是个格式化取值的宏，
% 本身永远非空。直接把原值抄过去，两个字段的精度声明一样，排出来也一样。
%    \begin{macrocode}
\AfterPreamble
  {
    \tl_if_empty:cT { hit@cover@date@zh@raw }
      { \cs_gset_eq:cc { hit@cover@date@zh@raw } { hit@date@zh@raw } }
  }
%</thesis-keylist>
%    \end{macrocode}
%
%
%
% \section{开题中期报告类的键清单（keylist）}
%
% \changes{v3.2a}{2026/08/10}{元信息字段的声明单独成文件。它说的是「有哪些字段」，
%   与 \file{src/setup/} 说的「字段取什么值」是一对，都是数据；
%   字段怎么存、旧名怎么兼容那套机制在 \file{src/utils/hit-meta.dtx}}
% 开题中期报告用到的元信息字段：字段名、旧名映射、\cs{hitsetup} 各族的路由，
% 以及摘要环境。都是清单，不含机制。
%
% ^^A ======================================================================
% ^^A Module report-keylist: 元信息字段（\hitsetup 接受的 key）（hit-report）
% ^^A ======================================================================
%    \begin{macrocode}
%<*report-keylist>
%% 中文封面

\hit@define@term{report-type-zh}
\hit@define@term{class-id}
\hit@define@date@term{enroll-date-zh}{zh}{auto}
%    \end{macrocode}
%
% 中文封面日期的精度在这里重设一遍。共享层那份按学位论文类的规矩写：本科到年月日、
% 研究生到年月。报告类里哈尔滨与威海本科到年月，深圳本科才到年月日，见 \issue{220}。
%
% \changes{v3.2a}{2026/08/16}{报告类封面日期的精度判断由 \texttt{harbin} 改成
%   \cs{hit@if@not@shenzhen}，威海本科的开题与中期跟着本部走}
% 精度这一段是取值时才求值的（见 \cs{hit@define@date@term}），所以这里可以写
% \cs{hit@if@not@shenzhen@TF}——它在 \file{src/setup/hit-report-options.dtx} 里
% 定义，按 \file{hithesis.ins} 的顺序晚于本文件，但排版时早就有了。
%    \begin{macrocode}
\hit@define@date@term{date-zh}{zh}
  {\hit@if@option@TF{bachelor}{\hit@if@not@shenzhen@TF{month}{day}}{month}}

%    \end{macrocode}
%
% \changes{v3.2a}{2026/08/10}{关键词的逗号拆分与 \cs{hitsetup} 的入口挪到
%   \file{src/utils/hit-meta.dtx}}
% 关键词怎么拆、\cs{hitsetup} 的占位定义都见 \file{src/utils/hit-meta.dtx}。
%
% \changes{v3.2a}{2026/08/10}{\cs{hitsetup} 改用 l3keys 做入口，加一层按职责
%   划分的子族路由，与 \file{hit-thesis.cls} 一致。子族暂时只转发，键仍由
%   \texttt{hit/class} 族处理}
% 入口换成 l3keys，根族路由到子族，子族一律把键转回
% \texttt{hit} 族。转发时值要重新加花括号，否则含逗号的值会被二次拆分成多个键。
%
% \changes{v3.2a}{2026/08/10}{元信息、可本地化字符串与参考文献选项正式注册进
%   对应的族，不再依赖子族兜底转发}
% 元信息字段与可本地化字符串。做法与 \file{hit-thesis.cls} 相同：
% 前者的处理器直接调用同名命令，后者的键名由宏名去掉 \texttt{hit@} 前缀、
% 其余 \texttt{@} 换成连字符得来。
%    \begin{macrocode}
\hit@declare@info@field
  {
    class-id , enroll-date-zh , report-type-zh
  }
\hit@declare@constant
  {
    report-header , report-front-header , shenzhen-master-school-bottom-mark ,
    bachelor-class-field-label , thesis-interim-check
  }
%    \end{macrocode}
%
% 参考文献的三个选项，改用连字符命名。
%
% \changes{v3.2a}{2026/08/16}{与学位论文类一样按 kebab-biblatex-gb7714 的词汇
%   重起名并接上 biblatex 一线的守卫}
% 键名、设值器与守卫跟学位论文类同一套，见那边的说明。报告类没有
% \option{compress-min} 与 \option{split-entry}，跟统一前一样。
%    \begin{macrocode}
\keys_define:nn { hit / bib }
  {
    numbering-align .code:n =
      {
        \str_if_eq:nnF {#1} { auto }
          {
            \hit@bib@key@guard:n { numbering-align }
            \str_gset:Nn \g_bib_labelalign_str {#1}
          }
      } ,
    uppercase-family .choices:nn = { true , false , auto }
      {
        \str_if_eq:VnF \l_keys_choice_tl { auto }
          {
            \hit@bib@key@guard:n { uppercase-family }
            \str_if_eq:VnTF \l_keys_choice_tl { true }
              { \bool_gset_true:N  \g_bib_authorupper_bool }
              { \bool_gset_false:N \g_bib_authorupper_bool }
          }
      } ,
    uppercase-family .default:n = true ,
    max-bib-names .code:n =
      {
        \str_if_eq:nnF {#1} { auto }
          {
            \hit@bib@key@guard:n { max-bib-names }
            \int_gset:Nn \g_bib_maxauthor_int {#1}
            \int_if_zero:nT { \g_bib_maxauthor_int }
              { \int_gdecr:N \g_bib_maxauthor_int }
          }
      } ,
    max-cite-names .code:n =
      {
        \str_if_eq:nnF {#1} { auto }
          {
            \hit@bib@key@guard:n { max-cite-names }
            \int_gset:Nn \g__hit_bib_maxcite_int {#1}
            \int_compare:nNnT \g__hit_bib_maxcite_int = 0
              { \int_gdecr:N \g__hit_bib_maxcite_int }
          }
      } ,
  }
%    \end{macrocode}
%
% \changes{v3.2a}{2026/08/11}{报告类把 \option{toc}/\option{english} 改成一条提示}
% 共享层的 \option{toc} 族三个键，报告类只做得了两个。一份报告只排一份目录，
% 语言跟着类选项 \option{lang} 走，没有「中文目录后面再跟一份英文的」这回事。
% 这里把 \option{english} 覆盖成一条提示，比让它静静地不起作用强。
%    \begin{macrocode}
\keys_define:nn { hit / toc }
  { english .code:n = { \msg_warning:nn { hithesis } { no-toc-english } } }
%    \end{macrocode}
%
% \changes{v3.2a}{2026/08/10}{加旧名到新名的映射表与废弃警告}
% 兼容期的处理，与 \file{hit-thesis.cls} 一致：旧名照常生效，只在日志里提示。
%    \begin{macrocode}
\prop_gset_from_keyval:Nn \g__hit_renamed_option_prop
  {
    raggedbottom   = { layout = { bottom = ragged~|~flush } } ,
    biblabelalign  = { bib = { numbering-align = ... } } ,
    bibauthorupper = { bib = { uppercase-family = true } } ,
    bibmaxauthor   = { bib = { max-bib-names = ... } } ,
  }
%    \end{macrocode}
%
% \changes{v3.2a}{2026/08/10}{元信息字段的旧名走废弃期：仍然可用，用一次提示一次}
% 旧名到新名的对照。旧名在 v3.1e 就是公开接口，写在用户的 \file{cover.tex} 里，
% 所以保留到 v4。
%    \begin{macrocode}
\hit@define@legacy@terms
  {
    ctype=report-type-zh ,
    csubject=major-zh ,
    cauthor=author-zh ,
    cstudentid=student-id ,
    cclassid=class-id ,
    caffil=affiliation-zh ,
    csupervisor=supervisor-zh ,
    cdate=date-zh ,
    cenrolldate=enroll-date-zh
  }
\hit@define@legacy@title{ctitlecover}{zh}{first}
\hit@define@legacy@title{ctitleone}{zh}{first}
\hit@define@legacy@title{ctitletwo}{zh}{second}
\prop_gset_from_keyval:Nn \g__hit_renamed_load_option_prop
  {
    font        = { font-en = ... } ,
    cjk-font    = { font-zh = ... } ,
    type        = { degree-level } ,
    newgeometry = { v2-layout~=~{~bachelor~,~graduate~=~one|two~} } ,
  }
\__hit_check_class_options:
%    \end{macrocode}
%
% 根族路由。\option{toc} 既是族名又是老的布尔选项名，所以它单独处理：
% 值是 \texttt{true}/\texttt{false} 时当老写法，其余当子族的键值表。
%    \begin{macrocode}
\keys_define:nn { hit }
  {
    info  .code:n = { \hit@keys@set:nn { hit / info } {#1} } ,
    style .code:n = { \hit@keys@set:nn { hit / style } {#1} } ,
    layout .code:n = { \hit@layout@set:n {#1} } ,
    format .code:n = { \hit@keys@set:nn { hit / format } {#1} } ,
    bib   .code:n =
      { \hit@keys@set:nn { hit / bib } {#1} } ,
    const .code:n = { \hit@keys@set:nn { hit / const } {#1} } ,
    struct .code:n =
      { \__hit_structure_reset_written: \hit@keys@set:nn { hit / struct } {#1} } ,
    thm .code:n = { \hit@theorem@style {#1} } ,
    toc .code:n =
      {
        \str_case:nnF {#1}
          {
            { true }  { \tl_gset:cn { g__hit_structure_contents_tl } { true } }
            { false } { \tl_gset:cn { g__hit_structure_contents_tl } { false } }
          }
          { \hit@keys@set:nn { hit / toc } {#1} }
      } ,
    toc .default:n = true ,
    unknown .code:n = { \__hit_legacy_key:n {#1} } ,
  }
\RenewDocumentCommand \hitsetup { O{} +m }
  {
    \hit@if@condition:nT {#1}
      { \hit@keys@set:nn { hit } {#2} \hit@bib@backend@apply: }
  }
%</report-keylist>
%    \end{macrocode}
%
%
% \Finale
